%=============================================================================!
% 12.7  HOW MUCH THE TEMPERATURE FLUCTUATES                                   !
%=============================================================================!
\section{How much the temperature fluctuates}
\label{sec:sm-kinetic}

The kinetic temperature of a run is never still. This section derives the
distribution of the kinetic energy in the canonical ensemble, finds that
a run at fixed energy fluctuates less, and turns the difference into a
measurement of the heat capacity. Chapter~13 tests every thermostat
against the distribution found here.

\subsection*{The kinetic energy in the canonical ensemble}

In the $N_{\mathrm f}$ free components $c_k$ of
\cref{sec:sm-equipartition}, the kinetic energy is $K =
\frac12\sum_kc_k^2$. The velocities with kinetic energy below $K$ fill
the inside of a sphere of radius $\sqrt{2K}$ in $N_{\mathrm f}$
dimensions, whose volume, as for the ideal gas of \cref{sec:sm-micro}, is
a constant times $K^{N_{\mathrm f}/2}$, and those between $K$ and $K +
\dd K$ fill a shell whose volume is the derivative of that volume times
$\dd K$, a constant times $K^{N_{\mathrm f}/2 - 1}\,\dd K$. Each
velocity in the shell has the Boltzmann factor $\mathrm{e}^{-\beta K}$,
so
\begin{equation}
   \mathcal{P}(K) \propto K^{N_{\mathrm f}/2 - 1}\,\mathrm{e}^{-K/\kB T} .
   \label{eq:sm-gamma}
\end{equation}
This is the gamma distribution. For one atom, $N_{\mathrm f} = 3$, it is
the $\sqrt{K}\,\mathrm{e}^{-K/\kB T}$ of \cref{eq:sm-atom-energy}; for an
atom moving over a surface, $N_{\mathrm f} = 2$, it is
$\mathrm{e}^{-K/\kB T}/\kB T$, and the probability that $K$ is at least
some $U_{\mathrm b}$ is $\mathrm{e}^{-U_{\mathrm b}/\kB T}$. Its mean and
variance follow as in \cref{sec:sm-canonical}, applied to the kinetic
energy alone, which is independent of the positions: with $a =
N_{\mathrm f}/2$, the factor $K^{a-1}$ that counts the velocities does not depend on $\beta$, so each
derivative of the normalising integral with respect to $\beta$ brings
down a factor $-K$, as each derivative of $Z$ brought down $-E_s$. The
normalising integral is, with $s = \beta K$,
\begin{equation*}
   Z_K = \int_0^\infty K^{a - 1}\mathrm{e}^{-\beta K}\,\dd K
   = \beta^{-a}\int_0^\infty s^{a - 1}\mathrm{e}^{-s}\,\dd s ,
\end{equation*}
and the last integral is a number, the gamma function of $a$, written
$\Gamma(a)$ (no relation to the phase point $\Gamma$), which
the function \texttt{kinetic\_\allowbreak energy\_\allowbreak density}
evaluates through its logarithm. So $\ln Z_K = -a\ln\beta + \text{a
constant}$, and
\begin{equation*}
   \ens{K} = -\frac{\dd\ln Z_K}{\dd\beta} = \frac a\beta
   = \tfrac12N_{\mathrm f}\kB T , \qquad
   \ens{K^2} - \ens{K}^2 = \frac{\dd^2\ln Z_K}{\dd\beta^2} = \frac a{\beta^2}
   = \tfrac12N_{\mathrm f}(\kB T)^2 .
\end{equation*}
The first is equipartition again. The second gives the relative
fluctuation of the kinetic energy, and of the kinetic temperature, which
is proportional to it:
\begin{equation}
   \frac{\sigma_K}{\ens{K}} = \sqrt{\frac{2}{N_{\mathrm f}}} .
   \label{eq:sm-relative}
\end{equation}
For one atom it is 0.82; for 256 argon atoms, with $N_{\mathrm f} = 765$,
0.0511. A canonical run of 256 atoms at \SI{135}{\kelvin} has a kinetic
temperature with a standard deviation of \SI{6.9}{\kelvin}, and the larger
the system, the smaller the spread.

\subsection*{At fixed energy, less}

In a run at fixed energy $E$, the kinetic and the potential energy share
$E$, as the two gases of \cref{sec:sm-micro} did. The probability that the
kinetic energy is $K$ is proportional to the number of velocities with
that kinetic energy, $K^{a - 1}$, times $\Omega_U(E - K)$, the volume of
the positions with potential energy $E - K$, which counts them as
$\Omega$ counted states in \cref{sec:sm-micro}. Its width follows as in
\cref{sec:sm-micro}, from the second derivative of its logarithm at the
peak, which is the sum of two terms, one from each partner:
\begin{align*}
   \frac{\dd^2}{\dd K^2}\bigl[(a - 1)\ln K\bigr] &= -\frac{a - 1}{K^2}
   \approx -\frac{a}{(a\kB T)^2} = -\frac{1}{\kB T^2C_K} ,
   && C_K = a\kB = \tfrac12N_{\mathrm f}\kB ,\\
   \frac{\dd^2}{\dd K^2}\ln\Omega_U(E - K)
   &= \frac{\dd^2\ln\Omega_U}{\dd U^2}
   = \frac{\dd}{\dd U}\Bigl(\frac{1}{\kB T}\Bigr)
   = -\frac{1}{\kB T^2C_U} ,
   && C_U = \frac{\dd U}{\dd T} ,
\end{align*}
using in the first $a - 1 \approx a$ for many freedoms and $K = a\kB T$ at
the peak, where, as for the two gases, both partners have the same
temperature. In the second, $U = E - K$, so each derivative with respect
to $K$ is minus that with respect to $U$ and the two signs cancel; then
come the definition of temperature, \cref{eq:sm-temperature}, applied to
the arrangements, and $\dd(1/\kB T)/\dd U = -(1/\kB T^2)\,\dd T/\dd U$ as
in \cref{sec:sm-canonical}. $C_K$ is the kinetic energy taken up per
degree and $C_U$ the potential energy, and together they make the heat
capacity, $C_V = C_K + C_U$. The logarithm of the Gaussian of
\cref{eq:sm-gauss} has the second derivative $-1/\sigma_x^2$, so the variance of $K$, written
$\ens{\delta K^2}_E$ with $\delta K = K - \ens{K}$, is minus the
reciprocal of the sum of the two curvatures:
\begin{equation}
   \ens{\delta K^2}_E = \frac{\kB T^2}{1/C_K + 1/C_U}
   = \kB T^2C_K\Bigl(1 - \frac{C_K}{C_V}\Bigr)
   = \tfrac12N_{\mathrm f}(\kB T)^2\Bigl(1 - \frac{N_{\mathrm f}\kB}
   {2C_V}\Bigr) ,
   \label{eq:sm-lpv}
\end{equation}
the result of Lebowitz, Percus and Verlet\cite{Lebowitz1967}. The bracket
is the factor by which the variance at fixed energy falls below the
canonical one, and its square root the factor for the standard
deviation. An ideal gas, with no potential energy, has $C_U = 0$
and a kinetic energy that cannot fluctuate at all; a harmonic crystal,
with $C_U = C_K$, fluctuates by $1/\sqrt2$ of the canonical amount.

\Cref{fig:sm-kinetic}(a) shows the kinetic energy of the 256 atoms of
liquid argon over \SI{50}{\pico\second}: its relative standard deviation
is 0.03068, against 0.05113 for the canonical ensemble, the ratio 0.600.
The heat capacity measured below, $C_V = 2.376\,N\kB$, with $C_K =
\frac12N_{\mathrm f}\kB = (765/512)\,N\kB = 1.4941\,N\kB$, predicts
$\sqrt{1 - 1.4941/2.376} = 0.609$.
The ratio stays near 0.6 for all four sizes of
\cref{fig:sm-kinetic}(b), from 0.564 to 0.600, while each spread falls
as $1/\sqrt N$.

\begin{figure}[htb]
\centering
\includegraphics{ch12_ensembles/kinetic}
\caption{The kinetic energy at fixed total energy. (a) Liquid argon, 256
atoms over \SI{50}{\pico\second}: the histogram of the kinetic energy
divided by its mean, against the canonical density of
\cref{eq:sm-gamma} with $N_{\mathrm f} = 765$ and the same mean (dashed).
(b) The relative standard deviation of the kinetic energy for 108, 256,
500 and 864 atoms (circles), against the canonical $\sqrt{2/N_{\mathrm
f}}$ (dashed).}
\label{fig:sm-kinetic}
\end{figure}

\subsection*{The heat capacity, two ways}

\Cref{eq:sm-lpv} solved for $C_V$,
\begin{equation*}
   C_V = \frac{\frac12N_{\mathrm f}\kB}{1 - 2\ens{\delta K^2}_E/
   \bigl(N_{\mathrm f}(\kB T)^2\bigr)} ,
\end{equation*}
gives the heat capacity from one run at fixed energy;
\texttt{statmech.heat\_capacity\_nve} computes it, with $T$ the run's
kinetic temperature. \Cref{fig:sm-heat} checks it against the definition,
$C_V = \dd E/\dd T$. Five runs of the 256 atoms, prepared near 100, 115,
130, 145 and \SI{160}{\kelvin} and run for \SI{50}{\pico\second} each,
give five pairs of mean energy and mean kinetic temperature, from
\SI{99.25}{\kelvin} to \SI{160.28}{\kelvin}. They lie on a straight line
to within \SI{0.029}{\electronvolt}, whose slope is
\SI{0.05242}{\electronvolt\per\kelvin}, or $C_V = 2.376\,N\kB$. The
spreads of the five runs give between 2.249 and $2.444\,N\kB$, with the
mean 2.347, and the four sizes near \SI{130}{\kelvin} between 2.198 and
2.333. The liquid lies between an ideal gas, with $\frac32N\kB$ from its
kinetic energy alone, and a harmonic crystal, with $3N\kB$ from equal
kinetic and potential shares. The single-run estimates scatter about the
slope by up to 7.5\%; whether that scatter is the noise of runs of
\SI{50}{\pico\second} is a question Chapter~15 answers with error bars.
In the canonical ensemble the total energy itself fluctuates, and
\cref{eq:sm-fluctuation} gives $C_V$ directly, which Chapter~13 measures
under its thermostats.

\begin{figure}[htb]
\centering
\includegraphics{ch12_ensembles/heat}
\caption{The heat capacity of liquid argon, 256 atoms. (a) The mean
energy against the mean kinetic temperature of five runs at fixed energy,
with the straight line through them, whose slope is $C_V$ (dashed).
(b) $C_V$ per atom, in units of $\kB$, from the spread of the kinetic
energy by \cref{eq:sm-lpv}, for the five runs (circles) and for runs of
108 to 864 atoms near \SI{130}{\kelvin} (squares), against the slope of
(a) (dashed).}
\label{fig:sm-heat}
\end{figure}

\begin{keyidea}
In the canonical ensemble the kinetic energy has the gamma distribution,
$\mathcal{P}(K) \propto K^{N_{\mathrm f}/2 - 1}\mathrm{e}^{-K/\kB T}$, with
mean $\frac12N_{\mathrm f}\kB T$ and relative spread
$\sqrt{2/N_{\mathrm f}}$. At fixed energy the kinetic and potential
energies share a fixed total, and the standard deviation of the kinetic
energy is $\sqrt{1 - N_{\mathrm f}\kB/(2C_V)}$ times the canonical one,
which gives the heat capacity from a single run.
\end{keyidea}

\notebook{12}{set the size and compare the kinetic energy at fixed energy with the canonical gamma distribution}

\begin{selfcheck}
By how much, relatively, does the kinetic temperature of 1000 atoms with
their drift removed fluctuate in the canonical ensemble?
\end{selfcheck}
